\documentclass[10pt,a4paper]{amsart}
\usepackage[margin=19mm]{geometry}
\usepackage{amsmath,amssymb,mathtools}
\usepackage[scr=rsfs,cal=euler]{mathalfa}
\usepackage{xfrac}
\usepackage[unicode,hidelinks,bookmarks=false]{hyperref}
\newcommand{\ie}{i.e.\ }
\newcommand{\cf}{cf.\ }
\newcommand{\resp}{respectively\ }
\newcommand{\Subsection}{\subsection}
\providecommand{\nobreakdash}{\nobreak\hbox{-}}
\setlength{\emergencystretch}{2em}
\multlinegap0pt
\usepackage{mathtools}
\usepackage{amssymb}
\usepackage{tikz}
\newtheorem{proposition}{Proposition}
\newtheorem{theorem}{Theorem}
\newcommand{\K}{\mathbb{K}}
\title{\bf Constraining Conformal Correlators}
\author{Viktoriia Borovik, Claire de Korte, Nathan Meurrens, Dmitrii Pavlov}
\date{Source: arXiv:2605.31491v1 (CC BY 4.0)}
\begin{document}
\maketitle

\noindent\emph{Source and licence.} \bf Constraining Conformal Correlators. Version arXiv:2605.31491v1 (CC BY 4.0), distributed by arXiv under the Creative Commons Attribution 4.0 International licence, http://creativecommons.org/licenses/by/4.0/, as recorded in the arXiv licence metadata for this exact version and shown on the abstract page https://arxiv.org/abs/2605.31491v1. Full licence notice: \texttt{licenses/arxiv-2605.31491-CC-BY-4.0.txt}.\par\smallskip

\noindent\textit{Editorial scope.} Counted unit: the article's theorem main (source0.tex lines 1418--1430). The article's complete original proof of it is delivered with it (source0.tex lines 1431--1477, one \texttt{proof} environment of 47 lines). No mathematics is written, repaired or re-derived: the statement and the proof are the article's own lines, in the article's own notation, and no retained line is substituted or re-flowed.  Retained uncounted as the source-local context the proof needs: lines 1257--1264; lines 1304--1311; lines 1377--1387; lines 1404--1412.  Nothing was omitted from any retained span, so the package carries no omission record.  No retained line contains a citation, so the wrapper carries no bibliography and no bibliography entry.  No retained line contains a figure environment, an image-inclusion command or drawing code, so no image is delivered and the package declares no asset dependency.  Editorial formatting only, in addition: an AMS \texttt{amsart} preamble that restates the article's own declarations and macros used by the retained lines, namely the environments \texttt{proposition}, \texttt{theorem} on separate counters, together with the standard packages those lines need.  The retained environments print in this document as Proposition 1, Theorem 1 in order of appearance, whereas the article prints the same items with the numbering of its own full document.  The complete native source of the article as distributed in the arXiv e-print for this exact version is bundled unchanged as \texttt{sources/201724/source0.tex}.
\begin{equation}\label{eq:transact}
 \forall \,  \alpha = (\alpha_1,\ldots,\alpha_n) \in \mathbb{G}_a^n\colon
  \quad
 \alpha \cdot P_i = P_i,
  \quad
  \alpha \cdot Z_i = Z_i + \alpha_i\,P_i,
  \quad i=1,\ldots,n.
\end{equation}

\begin{proposition}\label{prop:transformation-rules}
For $\alpha \in \mathbb{G}_a^n$, the generators of $S$ transform as follows:
\begin{equation}\label{eq:transform}
\alpha \cdot p_{ij} = p_{ij}, \qquad
\alpha \cdot q_{ij} = q_{ij} + \alpha_j\, p_{ij}, \qquad
\alpha \cdot r_{ij} = r_{ij} + \alpha_i\, q_{ij} + \alpha_j\, q_{ji} + \alpha_i \alpha_j\, p_{ij}.
\end{equation}
\end{proposition}

\begin{theorem}[Rosenlicht, 1956] \label{thm: Rosenlicht}
Let $G$ be an algebraic group acting on an irreducible variety~$X$.
There exists a finite set of rational invariants $M = \{f_1,\ldots,f_r\}$ that
separates orbits in general position. Moreover, the field of invariants
$\K(X)^G$ is generated by $M$, and
\begin{equation}\label{eq: trdegInvField}
\mathrm{trdeg}_\K\bigl(\K(X)^G\bigr)
  =
  \dim X - \max_{x\in X}\, \dim(G\cdot x).    
\end{equation}
\end{theorem}

\begin{proposition}\label{prop: inv}
The $n(2n-3)$  functions $P_{ij}$, $H_{ij}$, $\mathcal{V}_{ij}$ are $G$-invariant:
\[
  f(\alpha\cdot x) = f(x)
  \qquad
  \forall\,\alpha\in \mathbb{G}_a^n,\quad x\in \mathrm{Spec}(\mathcal{R}),\quad
  f\in\{P_{ij},\,H_{ij},\,\mathcal{V}_{ij}\}.
\]
\end{proposition}

\begin{theorem}\label{thm: main}
The field of rational invariants of the action induced by
\eqref{eq:transact} on $\mathrm{Spec}(\mathcal{R})$ is
\[
  \mathcal{F}^{G}
  \;=\;
  \K\Bigl(
    P_{ij}\;\,(i<j),\;\,
    H_{ij}\;\,(i<j),\;\,
    \mathcal{V}_{ij}\;\,(i,j\in[n],\;j\neq i,i+1)
  \Bigr).
\]
\end{theorem}

\begin{proof}
Invariance of the generators was established in
Proposition~\ref{prop: inv}, so the right‑hand side is contained in
$\mathcal{F}^G$. It therefore suffices to show that these invariants separate
general $G$‑orbits. In other words, to show that for any $x,x' \in \mathrm{Spec}(\mathcal R)$ and all admissible indices
\begin{equation}\label{eq: orbitseparators}
   P_{ij}(x')=P_{ij}(x),
  \quad
  H_{ij}(x')=H_{ij}(x),
  \quad
  \mathcal{V}_{ij}(x')=\mathcal{V}_{ij}(x),  
\end{equation}
if and only if  $x'=\alpha\cdot x$ for a unique $\alpha\in \mathbb{G}_a^n$. 
We consider equations~\eqref{eq: orbitseparators} on the dense open subset of $\mathrm{Spec}(\mathcal R)$ where the Gram matrix has maximal rank and the denominators of $\mathcal V_{ij}$ do not vanish.
On this
locus, the relations \eqref{eq: orbitseparators} are not implied by the Gram ideal.
In addition, since $P_{ij}=p_{ij}$, the equality $P_{ij}(x)=P_{ij}(x')$
implies that
\begin{equation}\label{eqn: pij equal}
p'_{ij}=p_{ij} \quad \text{for all} \quad  \;  i\neq j.
\end{equation}
Assuming \eqref{eqn: pij equal}, the equality $\mathcal V_{ij}(x')=\mathcal V_{ij}(x)$ then yields
$
  p_{ij}(q'_{i+1,i}-q_{i+1,i})
  =
  p_{i,i+1}(q'_{ji}-q_{ji})
$ for $j\neq i,i+1$.
That is, the vector
$(q'_{ji}-q_{ji})_{j\neq i}$ is proportional to $(p_{ij})_{j\neq i}$. Setting
\begin{equation*}\label{eq:alphai}
  \alpha_i := \frac{q'_{i+1, i}-q_{i+1, i}}{p_{i,i+1}},
\end{equation*}
it follows that
$
 q'_{ji}=q_{ji}+\alpha_i\,p_{ij}
$. Finally, substituting these expressions into the relations
$H_{ij}(x')=H_{ij}(x)$ gives
$$
  r'_{ij}
  =
  r_{ij}
  + \alpha_i q_{ij}
  + \alpha_j q_{ji}
  + \alpha_i\alpha_j p_{ij},
$$
Thus, by Proposition \ref{prop:transformation-rules} we have $x'=\alpha\cdot x$ for a unique $\alpha\in\mathbb G_a^n$, and the building blocks $P_{ij}, H_{ij}$ and $\mathcal{V}_{ij}$ separate general orbits. The claim now follows from Theorem~\ref{thm: Rosenlicht}.
\end{proof}

\end{document}
